\documentclass[11pt]{article}
\usepackage{amsmath, amssymb, geometry, hyperref}
\geometry{a4paper, margin=1in}

\title{Problem Sets: Null Spaces, Column Spaces, Row Spaces, and Linear Transformations}
\author{}
\date{}

\begin{document}

\maketitle
\tableofcontents
\newpage

\section{Concept Checklist \& Question Mapping}

This section outlines all the core concepts, problem types, and tricks found in the source material, along with the specific questions in this document that test each concept.

\begin{itemize}
    \item \textbf{Concept 1: Checking Membership in the Null Space ($Nul A$)}
    \begin{itemize}
        \item \textit{Description:} Verifying if a vector $w$ is in $Nul A$ by computing the matrix-vector product $Aw$ and checking if it equals the zero vector.
        \item \textit{Relevant Questions:} 1, 2, 3
    \end{itemize}
    
    \item \textbf{Concept 2: Checking Membership in the Column Space ($Col A$)}
    \begin{itemize}
        \item \textit{Description:} Verifying if a vector $w$ is in $Col A$ by setting up an augmented matrix $[A | w]$ and checking for consistency via row reduction.
        \item \textit{Relevant Questions:} 4, 5, 6
    \end{itemize}

    \item \textbf{Concept 3: Finding an Explicit Description (Basis) of the Null Space}
    \begin{itemize}
        \item \textit{Description:} Solving the homogeneous equation $Ax = 0$ by row reducing $A$ to RREF, identifying free variables, and writing the solution in parametric vector form.
        \item \textit{Relevant Questions:} 7, 8, 9, 10
    \end{itemize}

    \item \textbf{Concept 4: Finding a Basis for the Column Space}
    \begin{itemize}
        \item \textit{Description:} Row reducing $A$ to identify pivot columns, then selecting the corresponding columns from the \textit{original} matrix $A$.
        \item \textit{Relevant Questions:} 11, 12, 13
    \end{itemize}

    \item \textbf{Concept 5: Finding a Basis for the Row Space ($Row A$)}
    \begin{itemize}
        \item \textit{Description:} Row reducing $A$ to echelon form and taking the non-zero rows as the basis.
        \item \textit{Relevant Questions:} 14, 15, 16
    \end{itemize}

    \item \textbf{Concept 6: Proving/Disproving a Set is a Vector Space}
    \begin{itemize}
        \item \textit{Description:} Using the "Zero Vector Quick Check" or checking closure under addition and scalar multiplication.
        \item \textit{Relevant Questions:} 17, 18, 19
    \end{itemize}

    \item \textbf{Concept 7: Finding a Matrix $A$ such that a Set is $Col A$}
    \begin{itemize}
        \item \textit{Description:} Decomposing a generic vector defined by parameters into a linear combination of constant vectors, which become the columns of $A$.
        \item \textit{Relevant Questions:} 20, 21, 22
    \end{itemize}

    \item \textbf{Concept 8: Determining Subspace Dimensions ($k$)}
    \begin{itemize}
        \item \textit{Description:} Identifying that for an $m \times n$ matrix, $Nul A$ is a subspace of $\mathbb{R}^n$ (domain) and $Col A$ is a subspace of $\mathbb{R}^m$ (codomain).
        \item \textit{Relevant Questions:} 23, 24, 25
    \end{itemize}

    \item \textbf{Concept 9: Abstract Linear Transformations (Kernel and Range)}
    \begin{itemize}
        \item \textit{Description:} Applying the definitions of Kernel ($T(x) = 0$) and Range to abstract spaces like polynomials ($\mathbb{P}_n$) and matrices ($M_{m \times n}$).
        \item \textit{Relevant Questions:} 26, 27, 28
    \end{itemize}

    \item \textbf{Concept 10: Identifying Flaws and Conceptual Pitfalls}
    \begin{itemize}
        \item \textit{Description:} Diagnosing common mistakes, such as using RREF columns for $Col A$ or setting up an augmented matrix to check $Nul A$.
        \item \textit{Relevant Questions:} 29, 30
    \end{itemize}
\end{itemize}

\newpage

\section{Problem Sets}

\textbf{Type 1: Checking Membership in the Null Space}

\noindent \textbf{Problem 1:} Let $A = \begin{bmatrix} 1 & 2 \\ -1 & -2 \end{bmatrix}$. Determine if $w = \begin{bmatrix} 2 \\ -1 \end{bmatrix}$ is in $Nul A$.

\noindent \textbf{Problem 2:} Let $A = \begin{bmatrix} 2 & 0 & -4 \\ 1 & 1 & -1 \end{bmatrix}$. Determine if $w = \begin{bmatrix} 2 \\ 1 \\ 1 \end{bmatrix}$ is in $Nul A$.

\noindent \textbf{Problem 3:} Let $A = \begin{bmatrix} 1 & -1 & 0 \\ 0 & 2 & -2 \\ -1 & 0 & 1 \end{bmatrix}$. Determine if $w = \begin{bmatrix} 1 \\ 1 \\ 1 \end{bmatrix}$ is in $Nul A$.

\vspace{0.5cm}
\noindent \textbf{Type 2: Checking Membership in the Column Space}

\noindent \textbf{Problem 4:} Let $A = \begin{bmatrix} 1 & 2 \\ 3 & 6 \end{bmatrix}$. Determine if $w = \begin{bmatrix} 2 \\ 6 \end{bmatrix}$ is in $Col A$.

\noindent \textbf{Problem 5:} Let $A = \begin{bmatrix} 1 & 0 \\ 0 & 1 \\ 1 & 1 \end{bmatrix}$. Determine if $w = \begin{bmatrix} 1 \\ 2 \\ 4 \end{bmatrix}$ is in $Col A$.

\noindent \textbf{Problem 6:} Let $A = \begin{bmatrix} 2 & -1 & 3 \\ 4 & -2 & 6 \end{bmatrix}$. Determine if $w = \begin{bmatrix} 1 \\ 2 \end{bmatrix}$ is in $Col A$.

\vspace{0.5cm}
\noindent \textbf{Type 3: Finding an Explicit Description of the Null Space}

\noindent \textbf{Problem 7:} Find an explicit description of $Nul A$ by listing vectors that span the null space, where $A = \begin{bmatrix} 1 & -2 & 0 \\ 0 & 0 & 1 \end{bmatrix}$.

\noindent \textbf{Problem 8:} Find a basis for the null space of the $1 \times 4$ matrix $A = \begin{bmatrix} 1 & 2 & 3 & 4 \end{bmatrix}$.

\noindent \textbf{Problem 9:} Find an explicit description of $Nul A$ for $A = \begin{bmatrix} 1 & 0 & -1 \\ 2 & 0 & -2 \\ 3 & 0 & -3 \end{bmatrix}$.

\noindent \textbf{Problem 10:} Find a basis for $Nul A$ where $A = \begin{bmatrix} 2 & 4 & -2 \\ 1 & 2 & -1 \end{bmatrix}$.

\vspace{0.5cm}
\noindent \textbf{Type 4: Finding a Basis for the Column Space}

\noindent \textbf{Problem 11:} Find a basis for $Col A$ where $A = \begin{bmatrix} 1 & 2 & 3 \\ 0 & 1 & 2 \\ 0 & 0 & 0 \end{bmatrix}$.

\noindent \textbf{Problem 12:} Find a basis for $Col A$ where $A = \begin{bmatrix} 1 & -1 & 0 \\ -2 & 2 & 0 \\ 3 & -3 & 1 \end{bmatrix}$.

\noindent \textbf{Problem 13:} Find a basis for $Col A$ where $A = \begin{bmatrix} 2 & 4 \\ 1 & 2 \\ 3 & 6 \end{bmatrix}$.

\vspace{0.5cm}
\noindent \textbf{Type 5: Finding a Basis for the Row Space}

\noindent \textbf{Problem 14:} Find a basis for $Row A$ where $A = \begin{bmatrix} 1 & 2 & 3 \\ 0 & 1 & 2 \\ 0 & 0 & 0 \end{bmatrix}$.

\noindent \textbf{Problem 15:} Find a basis for $Row A$ where $A = \begin{bmatrix} 1 & -1 & 2 \\ -2 & 2 & -4 \end{bmatrix}$.

\noindent \textbf{Problem 16:} Find a basis for $Row A$ where $A = \begin{bmatrix} 0 & 1 & -1 \\ 2 & 4 & 0 \\ 1 & 2 & 0 \end{bmatrix}$.

\vspace{0.5cm}
\noindent \textbf{Type 6: Proving a Set is a Vector Space}

\noindent \textbf{Problem 17:} Determine if $W = \left\{ \begin{bmatrix} a \\ b \\ c \end{bmatrix} : a - 3b + c = 1 \right\}$ is a vector space.

\noindent \textbf{Problem 18:} Determine if $W = \left\{ \begin{bmatrix} 2s - t \\ s \\ 3t \end{bmatrix} : s, t \in \mathbb{R} \right\}$ is a vector space.

\noindent \textbf{Problem 19:} Determine if $W = \left\{ \begin{bmatrix} a \\ b \end{bmatrix} : a \ge 0, b \ge 0 \right\}$ is a vector space.

\vspace{0.5cm}
\noindent \textbf{Type 7: Finding a Matrix $A$ such that a Set is $Col A$}

\noindent \textbf{Problem 20:} Find a matrix $A$ such that the given set is $Col A$: $W = \left\{ \begin{bmatrix} 3x - y \\ x + 2y \\ -x \end{bmatrix} : x, y \in \mathbb{R} \right\}$.

\noindent \textbf{Problem 21:} Find a matrix $A$ such that the given set is $Col A$: $W = \left\{ \begin{bmatrix} a+b \\ a-b \\ 2a \\ 3b \end{bmatrix} : a, b \in \mathbb{R} \right\}$.

\noindent \textbf{Problem 22:} Find a matrix $A$ such that the given set is $Col A$: $W = \left\{ \begin{bmatrix} c_1 - 2c_2 + c_3 \\ c_2 - 4c_3 \end{bmatrix} : c_1, c_2, c_3 \in \mathbb{R} \right\}$.

\vspace{0.5cm}
\noindent \textbf{Type 8: Determining Subspace Dimensions}

\noindent \textbf{Problem 23:} Let $A$ be a $5 \times 3$ matrix. Find $k$ such that $Nul A$ is a subspace of $\mathbb{R}^k$, and find $j$ such that $Col A$ is a subspace of $\mathbb{R}^j$.

\noindent \textbf{Problem 24:} Let $A$ be a $1 \times 4$ matrix. Find $k$ such that $Nul A$ is a subspace of $\mathbb{R}^k$, and find $j$ such that $Col A$ is a subspace of $\mathbb{R}^j$.

\noindent \textbf{Problem 25:} Let $A$ be an $n \times n$ matrix. If $Nul A$ contains only the zero vector, what is the dimension of $Col A$?

\vspace{0.5cm}
\noindent \textbf{Type 9: Abstract Linear Transformations}

\noindent \textbf{Problem 26:} Define $T: \mathbb{P}_2 \to \mathbb{R}$ by $T(p) = p(1)$. Find a basis for the kernel of $T$.

\noindent \textbf{Problem 27:} Let $M_{2 \times 2}$ be the vector space of all $2 \times 2$ matrices. Define $T: M_{2 \times 2} \to M_{2 \times 2}$ by $T(A) = A - A^T$. Describe the kernel of $T$ and provide a basis for it.

\noindent \textbf{Problem 28:} Define a linear transformation $T: \mathbb{R}^3 \to \mathbb{R}^2$ by $T(x) = \begin{bmatrix} 1 & 0 & -1 \\ 0 & 1 & 1 \end{bmatrix} x$. Describe the range of $T$.

\vspace{0.5cm}
\noindent \textbf{Type 10: Advanced Diagnostic Testing ("Find the Flaw")}

\noindent \textbf{Problem 29:} \textit{Flawed Solution:} To find a basis for $Col A$ where $A = \begin{bmatrix} 1 & 3 \\ 2 & 6 \end{bmatrix}$, I row reduced $A$ to get $U = \begin{bmatrix} 1 & 3 \\ 0 & 0 \end{bmatrix}$. The non-zero columns of the reduced matrix form the basis, so the basis is $\begin{bmatrix} 1 \\ 0 \end{bmatrix}$. \\
\textbf{Task:} Find the error, explain why it is wrong in one sentence, and provide the correct solution.

\noindent \textbf{Problem 30:} \textit{Flawed Solution:} To check if $v = \begin{bmatrix} 1 \\ 1 \end{bmatrix}$ is in $Nul A$ for $A = \begin{bmatrix} 1 & -1 \\ 2 & -2 \end{bmatrix}$, we set up the augmented matrix $\begin{bmatrix} 1 & -1 & | & 1 \\ 2 & -2 & | & 1 \end{bmatrix}$. Row reducing gives $\begin{bmatrix} 1 & -1 & | & 1 \\ 0 & 0 & | & -1 \end{bmatrix}$. The system is inconsistent, so $v$ is not in $Nul A$. \\
\textbf{Task:} Find the error, explain why it is wrong in one sentence, and provide the correct solution.

\newpage

\section{Detailed Solutions}

\textbf{Solution 1:} To determine if $w \in Nul A$, we must check if $Aw = 0$.
$$ Aw = \begin{bmatrix} 1 & 2 \\ -1 & -2 \end{bmatrix} \begin{bmatrix} 2 \\ -1 \end{bmatrix} = \begin{bmatrix} 1(2) + 2(-1) \\ -1(2) + (-2)(-1) \end{bmatrix} = \begin{bmatrix} 2 - 2 \\ -2 + 2 \end{bmatrix} = \begin{bmatrix} 0 \\ 0 \end{bmatrix} $$
Since $Aw = 0$, the vector $w$ is in $Nul A$.

\vspace{0.3cm}
\textbf{Solution 2:} We check if $Aw = 0$.
$$ Aw = \begin{bmatrix} 2 & 0 & -4 \\ 1 & 1 & -1 \end{bmatrix} \begin{bmatrix} 2 \\ 1 \\ 1 \end{bmatrix} = \begin{bmatrix} 2(2) + 0(1) - 4(1) \\ 1(2) + 1(1) - 1(1) \end{bmatrix} = \begin{bmatrix} 4 - 4 \\ 2 + 1 - 1 \end{bmatrix} = \begin{bmatrix} 0 \\ 2 \end{bmatrix} $$
Since $Aw \neq 0$, the vector $w$ is \textbf{not} in $Nul A$.

\vspace{0.3cm}
\textbf{Solution 3:} We check if $Aw = 0$.
$$ Aw = \begin{bmatrix} 1 & -1 & 0 \\ 0 & 2 & -2 \\ -1 & 0 & 1 \end{bmatrix} \begin{bmatrix} 1 \\ 1 \\ 1 \end{bmatrix} = \begin{bmatrix} 1 - 1 + 0 \\ 0 + 2 - 2 \\ -1 + 0 + 1 \end{bmatrix} = \begin{bmatrix} 0 \\ 0 \\ 0 \end{bmatrix} $$
Since $Aw = 0$, the vector $w$ is in $Nul A$.

\vspace{0.3cm}
\textbf{Solution 4:} To check if $w \in Col A$, we determine if $Ax = w$ is consistent by setting up an augmented matrix:
$$ \begin{bmatrix} 1 & 2 & | & 2 \\ 3 & 6 & | & 6 \end{bmatrix} $$
Row reduce by replacing Row 2 with Row 2 - 3(Row 1):
$$ \begin{bmatrix} 1 & 2 & | & 2 \\ 0 & 0 & | & 0 \end{bmatrix} $$
The system is consistent (no row of the form $[0 \ 0 \ | \ c]$ where $c \neq 0$). Therefore, $w \in Col A$.

\vspace{0.3cm}
\textbf{Solution 5:} Set up the augmented matrix $[A | w]$:
$$ \begin{bmatrix} 1 & 0 & | & 1 \\ 0 & 1 & | & 2 \\ 1 & 1 & | & 4 \end{bmatrix} $$
Row reduce by replacing Row 3 with Row 3 - Row 1:
$$ \begin{bmatrix} 1 & 0 & | & 1 \\ 0 & 1 & | & 2 \\ 0 & 1 & | & 3 \end{bmatrix} $$
Replace Row 3 with Row 3 - Row 2:
$$ \begin{bmatrix} 1 & 0 & | & 1 \\ 0 & 1 & | & 2 \\ 0 & 0 & | & 1 \end{bmatrix} $$
The last row corresponds to the equation $0 = 1$, which is inconsistent. Therefore, $w$ is \textbf{not} in $Col A$.

\vspace{0.3cm}
\textbf{Solution 6:} Set up the augmented matrix $[A | w]$:
$$ \begin{bmatrix} 2 & -1 & 3 & | & 1 \\ 4 & -2 & 6 & | & 2 \end{bmatrix} $$
Replace Row 2 with Row 2 - 2(Row 1):
$$ \begin{bmatrix} 2 & -1 & 3 & | & 1 \\ 0 & 0 & 0 & | & 0 \end{bmatrix} $$
The system is consistent. Therefore, $w \in Col A$.

\vspace{0.3cm}
\textbf{Solution 7:} We solve $Ax = 0$. The matrix is already in RREF. The equations are:
$$ x_1 - 2x_2 = 0 \implies x_1 = 2x_2 $$
$$ x_3 = 0 $$
The free variable is $x_2$. Writing the solution in parametric vector form:
$$ x = \begin{bmatrix} x_1 \\ x_2 \\ x_3 \end{bmatrix} = \begin{bmatrix} 2x_2 \\ x_2 \\ 0 \end{bmatrix} = x_2 \begin{bmatrix} 2 \\ 1 \\ 0 \end{bmatrix} $$
Thus, $Nul A = \text{Span} \left\{ \begin{bmatrix} 2 \\ 1 \\ 0 \end{bmatrix} \right\}$.

\vspace{0.3cm}
\textbf{Solution 8:} We solve $Ax = 0$. The equation is $x_1 + 2x_2 + 3x_3 + 4x_4 = 0$.
Solving for the pivot variable $x_1$:
$$ x_1 = -2x_2 - 3x_3 - 4x_4 $$
The free variables are $x_2, x_3, x_4$. In parametric vector form:
$$ x = \begin{bmatrix} x_1 \\ x_2 \\ x_3 \\ x_4 \end{bmatrix} = \begin{bmatrix} -2x_2 - 3x_3 - 4x_4 \\ x_2 \\ x_3 \\ x_4 \end{bmatrix} = x_2 \begin{bmatrix} -2 \\ 1 \\ 0 \\ 0 \end{bmatrix} + x_3 \begin{bmatrix} -3 \\ 0 \\ 1 \\ 0 \end{bmatrix} + x_4 \begin{bmatrix} -4 \\ 0 \\ 0 \\ 1 \end{bmatrix} $$
A basis for $Nul A$ is $\left\{ \begin{bmatrix} -2 \\ 1 \\ 0 \\ 0 \end{bmatrix}, \begin{bmatrix} -3 \\ 0 \\ 1 \\ 0 \end{bmatrix}, \begin{bmatrix} -4 \\ 0 \\ 0 \\ 1 \end{bmatrix} \right\}$.

\vspace{0.3cm}
\textbf{Solution 9:} We solve $Ax = 0$ by row reducing $A$:
$$ \begin{bmatrix} 1 & 0 & -1 \\ 2 & 0 & -2 \\ 3 & 0 & -3 \end{bmatrix} \sim \begin{bmatrix} 1 & 0 & -1 \\ 0 & 0 & 0 \\ 0 & 0 & 0 \end{bmatrix} $$
The equation is $x_1 - x_3 = 0 \implies x_1 = x_3$. The free variables are $x_2$ and $x_3$.
$$ x = \begin{bmatrix} x_1 \\ x_2 \\ x_3 \end{bmatrix} = \begin{bmatrix} x_3 \\ x_2 \\ x_3 \end{bmatrix} = x_2 \begin{bmatrix} 0 \\ 1 \\ 0 \end{bmatrix} + x_3 \begin{bmatrix} 1 \\ 0 \\ 1 \end{bmatrix} $$
Thus, $Nul A = \text{Span} \left\{ \begin{bmatrix} 0 \\ 1 \\ 0 \end{bmatrix}, \begin{bmatrix} 1 \\ 0 \\ 1 \end{bmatrix} \right\}$.

\vspace{0.3cm}
\textbf{Solution 10:} Row reduce $A$ to solve $Ax = 0$:
$$ \begin{bmatrix} 2 & 4 & -2 \\ 1 & 2 & -1 \end{bmatrix} \sim \begin{bmatrix} 1 & 2 & -1 \\ 2 & 4 & -2 \end{bmatrix} \sim \begin{bmatrix} 1 & 2 & -1 \\ 0 & 0 & 0 \end{bmatrix} $$
The equation is $x_1 + 2x_2 - x_3 = 0 \implies x_1 = -2x_2 + x_3$. Free variables are $x_2, x_3$.
$$ x = x_2 \begin{bmatrix} -2 \\ 1 \\ 0 \end{bmatrix} + x_3 \begin{bmatrix} 1 \\ 0 \\ 1 \end{bmatrix} $$
A basis for $Nul A$ is $\left\{ \begin{bmatrix} -2 \\ 1 \\ 0 \end{bmatrix}, \begin{bmatrix} 1 \\ 0 \\ 1 \end{bmatrix} \right\}$.

\vspace{0.3cm}
\textbf{Solution 11:} The matrix $A = \begin{bmatrix} 1 & 2 & 3 \\ 0 & 1 & 2 \\ 0 & 0 & 0 \end{bmatrix}$ is already in echelon form. The pivot columns are columns 1 and 2. We take the corresponding columns from the original matrix $A$.
Basis for $Col A$: $\left\{ \begin{bmatrix} 1 \\ 0 \\ 0 \end{bmatrix}, \begin{bmatrix} 2 \\ 1 \\ 0 \end{bmatrix} \right\}$.

\vspace{0.3cm}
\textbf{Solution 12:} Row reduce $A$ to find the pivot columns:
$$ \begin{bmatrix} 1 & -1 & 0 \\ -2 & 2 & 0 \\ 3 & -3 & 1 \end{bmatrix} \sim \begin{bmatrix} 1 & -1 & 0 \\ 0 & 0 & 0 \\ 0 & 0 & 1 \end{bmatrix} \sim \begin{bmatrix} 1 & -1 & 0 \\ 0 & 0 & 1 \\ 0 & 0 & 0 \end{bmatrix} $$
The pivots are in columns 1 and 3. We must use the columns from the \textbf{original} matrix $A$.
Basis for $Col A$: $\left\{ \begin{bmatrix} 1 \\ -2 \\ 3 \end{bmatrix}, \begin{bmatrix} 0 \\ 0 \\ 1 \end{bmatrix} \right\}$.

\vspace{0.3cm}
\textbf{Solution 13:} Row reduce $A$:
$$ \begin{bmatrix} 2 & 4 \\ 1 & 2 \\ 3 & 6 \end{bmatrix} \sim \begin{bmatrix} 1 & 2 \\ 2 & 4 \\ 3 & 6 \end{bmatrix} \sim \begin{bmatrix} 1 & 2 \\ 0 & 0 \\ 0 & 0 \end{bmatrix} $$
The pivot is in column 1. Using the original column:
Basis for $Col A$: $\left\{ \begin{bmatrix} 2 \\ 1 \\ 3 \end{bmatrix} \right\}$.

\vspace{0.3cm}
\textbf{Solution 14:} The row space is the span of the rows. Row operations do not change the row space. The matrix $A = \begin{bmatrix} 1 & 2 & 3 \\ 0 & 1 & 2 \\ 0 & 0 & 0 \end{bmatrix}$ is in echelon form. The non-zero rows form the basis.
Basis for $Row A$: $\left\{ \begin{bmatrix} 1 \\ 2 \\ 3 \end{bmatrix}, \begin{bmatrix} 0 \\ 1 \\ 2 \end{bmatrix} \right\}$.

\vspace{0.3cm}
\textbf{Solution 15:} Row reduce $A$:
$$ \begin{bmatrix} 1 & -1 & 2 \\ -2 & 2 & -4 \end{bmatrix} \sim \begin{bmatrix} 1 & -1 & 2 \\ 0 & 0 & 0 \end{bmatrix} $$
The non-zero row forms the basis.
Basis for $Row A$: $\left\{ \begin{bmatrix} 1 \\ -1 \\ 2 \end{bmatrix} \right\}$.

\vspace{0.3cm}
\textbf{Solution 16:} Row reduce $A$:
$$ \begin{bmatrix} 0 & 1 & -1 \\ 2 & 4 & 0 \\ 1 & 2 & 0 \end{bmatrix} \sim \begin{bmatrix} 1 & 2 & 0 \\ 2 & 4 & 0 \\ 0 & 1 & -1 \end{bmatrix} \sim \begin{bmatrix} 1 & 2 & 0 \\ 0 & 0 & 0 \\ 0 & 1 & -1 \end{bmatrix} \sim \begin{bmatrix} 1 & 2 & 0 \\ 0 & 1 & -1 \\ 0 & 0 & 0 \end{bmatrix} $$
The non-zero rows of the echelon form are the basis.
Basis for $Row A$: $\left\{ \begin{bmatrix} 1 \\ 2 \\ 0 \end{bmatrix}, \begin{bmatrix} 0 \\ 1 \\ -1 \end{bmatrix} \right\}$.

\vspace{0.3cm}
\textbf{Solution 17:} Use the "Zero Vector Quick Check". For the zero vector, $a=0, b=0, c=0$. Substituting into the condition:
$$ 0 - 3(0) + 0 = 0 \neq 1 $$
Since the zero vector is not in $W$, $W$ is \textbf{not} a vector space.

\vspace{0.3cm}
\textbf{Solution 18:} We can decompose the generic vector in $W$:
$$ \begin{bmatrix} 2s - t \\ s \\ 3t \end{bmatrix} = s \begin{bmatrix} 2 \\ 1 \\ 0 \end{bmatrix} + t \begin{bmatrix} -1 \\ 0 \\ 3 \end{bmatrix} $$
Since $W$ is the span of two vectors, it is a subspace of $\mathbb{R}^3$. Therefore, $W$ \textbf{is} a vector space.

\vspace{0.3cm}
\textbf{Solution 19:} Check closure under scalar multiplication. Let $v = \begin{bmatrix} 1 \\ 1 \end{bmatrix} \in W$. Let $c = -1$.
$$ cv = -1 \begin{bmatrix} 1 \\ 1 \end{bmatrix} = \begin{bmatrix} -1 \\ -1 \end{bmatrix} $$
Since $-1 \not\ge 0$, $cv$ is not in $W$. $W$ is \textbf{not} a vector space.

\vspace{0.3cm}
\textbf{Solution 20:} Decompose the generic vector into a linear combination using parameters $x$ and $y$ as weights:
$$ \begin{bmatrix} 3x - y \\ x + 2y \\ -x \end{bmatrix} = x \begin{bmatrix} 3 \\ 1 \\ -1 \end{bmatrix} + y \begin{bmatrix} -1 \\ 2 \\ 0 \end{bmatrix} $$
Construct matrix $A$ using these vectors as columns:
$$ A = \begin{bmatrix} 3 & -1 \\ 1 & 2 \\ -1 & 0 \end{bmatrix} $$

\vspace{0.3cm}
\textbf{Solution 21:} Decompose the vector:
$$ \begin{bmatrix} a+b \\ a-b \\ 2a \\ 3b \end{bmatrix} = a \begin{bmatrix} 1 \\ 1 \\ 2 \\ 0 \end{bmatrix} + b \begin{bmatrix} 1 \\ -1 \\ 0 \\ 3 \end{bmatrix} $$
Construct matrix $A$:
$$ A = \begin{bmatrix} 1 & 1 \\ 1 & -1 \\ 2 & 0 \\ 0 & 3 \end{bmatrix} $$

\vspace{0.3cm}
\textbf{Solution 22:} Decompose the vector:
$$ \begin{bmatrix} c_1 - 2c_2 + c_3 \\ c_2 - 4c_3 \end{bmatrix} = c_1 \begin{bmatrix} 1 \\ 0 \end{bmatrix} + c_2 \begin{bmatrix} -2 \\ 1 \end{bmatrix} + c_3 \begin{bmatrix} 1 \\ -4 \end{bmatrix} $$
Construct matrix $A$:
$$ A = \begin{bmatrix} 1 & -2 & 1 \\ 0 & 1 & -4 \end{bmatrix} $$

\vspace{0.3cm}
\textbf{Solution 23:} The matrix $A$ has dimensions $5 \times 3$ (5 rows, 3 columns).
\begin{itemize}
    \item $Nul A$ consists of vectors $x$ such that $Ax = 0$. $x$ must have 3 entries to match the columns of $A$. Thus, $Nul A$ is a subspace of $\mathbb{R}^3$ ($k=3$).
    \item $Col A$ consists of output vectors $b$ such that $Ax = b$. $b$ will have 5 entries to match the rows of $A$. Thus, $Col A$ is a subspace of $\mathbb{R}^5$ ($j=5$).
\end{itemize}

\vspace{0.3cm}
\textbf{Solution 24:} The matrix $A$ has dimensions $1 \times 4$.
\begin{itemize}
    \item $x$ must have 4 entries. $Nul A$ is a subspace of $\mathbb{R}^4$ ($k=4$).
    \item $b$ must have 1 entry. $Col A$ is a subspace of $\mathbb{R}^1$ or $\mathbb{R}$ ($j=1$).
\end{itemize}

\vspace{0.3cm}
\textbf{Solution 25:} By the Rank-Nullity Theorem (Dimension Theorem), the dimension of the null space plus the dimension of the column space equals the total number of columns ($n$). If $Nul A = \{0\}$, its dimension is 0. Therefore, the dimension of $Col A$ must be $n$. (The column space is all of $\mathbb{R}^n$).

\vspace{0.3cm}
\textbf{Solution 26:} The kernel of $T$ is the set of all polynomials $p \in \mathbb{P}_2$ such that $T(p) = 0$. This means $p(1) = 0$.
A polynomial in $\mathbb{P}_2$ has the form $p(t) = at^2 + bt + c$.
$$ p(1) = a(1)^2 + b(1) + c = a + b + c = 0 \implies c = -a - b $$
Substitute $c$ back into the polynomial:
$$ p(t) = at^2 + bt + (-a - b) = a(t^2 - 1) + b(t - 1) $$
A basis for the kernel is $\{t^2 - 1, t - 1\}$.

\vspace{0.3cm}
\textbf{Solution 27:} The kernel of $T$ is the set of matrices $A$ such that $T(A) = 0$.
$$ A - A^T = 0 \implies A = A^T $$
This is the definition of a symmetric matrix. Thus, the kernel of $T$ is the set of all symmetric $2 \times 2$ matrices. A standard basis for this space is:
$$ \left\{ \begin{bmatrix} 1 & 0 \\ 0 & 0 \end{bmatrix}, \begin{bmatrix} 0 & 1 \\ 1 & 0 \end{bmatrix}, \begin{bmatrix} 0 & 0 \\ 0 & 1 \end{bmatrix} \right\} $$

\vspace{0.3cm}
\textbf{Solution 28:} The range of $T$ is exactly the column space of the standard matrix $A = \begin{bmatrix} 1 & 0 & -1 \\ 0 & 1 & 1 \end{bmatrix}$.
The matrix $A$ is already in RREF and has a pivot in every row (columns 1 and 2). Since $Col A$ spans $\mathbb{R}^2$, the range of $T$ is all of $\mathbb{R}^2$.

\vspace{0.3cm}
\textbf{Solution 29:} \textbf{Error:} The solution uses the columns of the reduced matrix $U$ instead of the original matrix $A$. \\
\textbf{Correct Solution:} Row operations change the column space, so you must use the pivot columns of the \textit{original} matrix. The pivot is in the first column, so the correct basis is the first column of $A$, which is $\begin{bmatrix} 1 \\ 2 \end{bmatrix}$.

\vspace{0.3cm}
\textbf{Solution 30:} \textbf{Error:} The solution tests if $v \in Col A$ by setting up an augmented matrix, instead of testing if $v \in Nul A$ by computing $Av$. \\
\textbf{Correct Solution:} To check if $v \in Nul A$, compute $Av$.
$$ Av = \begin{bmatrix} 1 & -1 \\ 2 & -2 \end{bmatrix} \begin{bmatrix} 1 \\ 1 \end{bmatrix} = \begin{bmatrix} 1 - 1 \\ 2 - 2 \end{bmatrix} = \begin{bmatrix} 0 \\ 0 \end{bmatrix} $$
Since $Av = 0$, $v$ is indeed in $Nul A$.

\end{document}